% ════════════════════════════════════════════
% 模型检验与分析（6.模型检验.tex）写作规则 —— 模板内嵌，AI 先读本文件再写
% ════════════════════════════════════════════
\section{模型检验与分析}

为验证本文构建的模型的有效性和稳健性，从误差分析与灵敏度分析两个维度进行检验：误差分析考察各模型在训练集与检验集上的拟合精度与泛化能力，灵敏度分析考察模型输出对关键输入与参数扰动的稳定程度。

\subsection{误差分析}
\label{sec:check}

对问题一的岭正则线性混合回归、问题二的经典/广义标度律、问题四的 Loss—Benchmark 桥接与面板回归，分别计算决定系数 $R^2$、均方根误差 RMSE、平均绝对误差 MAE 与平均绝对百分比误差 MAPE，结果如表\ref{tab:error}所示。

\begin{longtable}{>{\raggedright\arraybackslash}p{0.40\textwidth}>{\centering\arraybackslash}p{0.11\textwidth}>{\centering\arraybackslash}p{0.12\textwidth}>{\centering\arraybackslash}p{0.12\textwidth}>{\raggedright\arraybackslash}p{0.13\textwidth}}
  \caption{模型误差分析结果}
  \label{tab:error} \\
  \toprule
  模型（数据集） & $R^2$ / $r$ & RMSE & MAE & MAPE(\%) \\
  \midrule
  \endfirsthead
  \caption{模型误差分析结果（续）} \\
  \toprule
  模型（数据集） & $R^2$ / $r$ & RMSE & MAE & MAPE(\%) \\
  \midrule
  \endhead
  \bottomrule
  \endfoot
  问题一岭回归（训练 1M，$n=512$） & 0.629 & 0.551 & 0.403 & 8.54 \\
  问题一岭回归（检验 1M，$n=256$） & 0.619 & 0.523 & 0.386 & 8.43 \\
  问题一跨尺度 Pearson（60M / 1B） & 0.782 / 0.651 & — & — & — \\
  问题二经典标度律（B1，$n=1{,}176$） & 0.9999998 & $1.5\times10^{-4}$ & $1.2\times10^{-4}$ & 0.004 \\
  问题二经典标度律留出（4 小模型$\to$4 大模型） & 0.9999997 & $1.5\times10^{-4}$ & — & 0.005 \\
  问题二广义标度律（B1+B6 联合，$n=1{,}536$） & 0.990 & 0.036 & 0.024 & 0.66 \\
  问题二跨族验证（B4，$n=57$） & 0.975 & 0.293 & — & 8.83 \\
  问题二文献验证（B5，$n=44$） & 0.910 & 0.193 & — & 7.99 \\
  问题四 Loss—Benchmark 桥接（C6，$n=75$） & 0.26 & — & — & 67.79 \\
  问题四面板回归（C1 开源筛选，$n=2{,}346$） & 0.475 & — & — & — \\
  问题四面板回归（C3 长周期，$n=4{,}586$） & 0.424 & — & — & — \\
  问题四前沿 logistic 外推 & 0.975 & — & — & — \\
  问题三优化约束残差（$|C_{\text{total}}/C_{\text{budget}}-1|$） & — & $<10^{-6}$ & — & — \\
\end{longtable}

由表\ref{tab:error}可知：（1）问题一的岭回归在 13 个损失域上的平均训练 $R^2=0.629$、同尺度检验 $R^2=0.619$，训练与检验精度基本一致（差距 $0.010$），未出现明显过拟合；MAPE 约 $8.5\%$，相对误差可控。（2）问题二的经典标度律在 B1 上 $R^2\approx1.000$、MAPE $0.004\%$，且以小模型训练、大模型留出的验证结果几乎不变（$R^2=0.9999997$）——但如 5.2 节所述，B1 的拟合残差标准差仅 $1.5\times10^{-4}$ 而损失标准差为 $0.342$，说明该附件近乎无噪声，此处的极高精度只能证明函数形式自洽，真正的泛化证据来自 B4（跨族真实数据，$r=0.975$、MAPE $8.8\%$）与 B5（文献数据，$r=0.910$、MAPE $8.0\%$）。（3）引入质量项后的广义标度律在 $1{,}536$ 个联合样本上 $R^2=0.990$、MAPE $0.66\%$，说明加入质量维未牺牲拟合精度。（4）问题四的 Loss—Benchmark 桥接 $R^2$ 仅 $0.26$、MAPE 达 $67.8\%$，这是本文误差最大的环节，其成因已在 5.4 节说明（中可比样本的验证集不统一、高可比样本的得分变化幅度过小），因此桥接结果仅用于方向性判定。（5）问题三的优化模型在所有求解点上均满足预算约束（相对残差 $<10^{-6}$），无约束违反。

\subsection{灵敏度分析}

为考察模型结论对关键输入与假设的稳健性，分别对赋权方案、冲突阈值、正则系数、质量成本函数形式、质量基线口径、上下文长度与算力预算进行扰动分析，结果如表\ref{tab:sensitivity}所示。

\begin{longtable}{>{\raggedright\arraybackslash}p{0.30\textwidth}>{\raggedright\arraybackslash}p{0.36\textwidth}>{\raggedright\arraybackslash}p{0.26\textwidth}}
  \caption{灵敏度分析结果}
  \label{tab:sensitivity} \\
  \toprule
  参数/设定变化 & 输出变化 & 灵敏度评估 \\
  \midrule
  \endfirsthead
  \caption{灵敏度分析结果（续）} \\
  \toprule
  参数/设定变化 & 输出变化 & 灵敏度评估 \\
  \midrule
  \endhead
  \bottomrule
  \endfoot
  赋权方案：熵权法 $\to$ 等权平均 $\to$ PCA & 域级质量排序一致性 $0.679\to0.679$；样本级得分相关 $0.854/0.881$；得分均值 $0.495\to0.639\to0.520$ & 中等敏感：整体格局稳定，但领域名次会变化 \\
  指标方向：ad\_en 按“第 1 类=无广告”（正向）$\to$ 若误判为负向 & 领域质量排序由「arxiv $>$ stackexchange $>$ c4 $>$ commoncrawl $>$ book $\approx$ wikipedia $>$ github」变为「c4 $>$ commoncrawl $>$ arxiv $>$ stackexchange $>$ book $>$ wikipedia $>$ github」 & 高度敏感：方向判定必须用实证证据 \\
  冲突阈值 $\tau$：$0.5\to1.0\to1.5\to2.0\to2.5$ & 冲突率 $0.670\to0.402\to0.223\to0.121\to0.067$；与独立基准之比 $0.93\to0.84\to0.77\to0.77\to0.86$ & 结论稳健：任一阈值下冲突率均显著低于独立基准 \\
  岭正则系数 $\lambda$：$10^{-4}\sim10^{-1}$ & 训练 $R^2$ 由 $0.6291$ 降至 $0.6238$，检验 $R^2$ 由 $0.6187$ 降至 $0.6157$ & 不敏感，估计稳定 \\
  质量成本函数：指数型 $\to$ 幂函数型 $\to$ 对数渐进型（$C=10^{19}$） & $Q^{*}=0.660\to0.551\to1.000$，$L^{*}$ 差异 $<0.06$；质量份额 $12.9\%\to0.0\%\to31.6\%$ & 低预算档较敏感 \\
  同上，但 $C\ge10^{22}$ & 三种函数下 $Q^{*}$ 均为 $1.000$，$N^{*}$、$D^{*}$ 差异 $<10\%$ & 不敏感，结论稳健 \\
  质量基线口径：配比加权 $Q_0=0.5510\to$ 样本均值 $0.4952$ & 三档预算下 $Q^{*}$：$0.551/1.000/1.000\to0.504/1.000/1.000$，$N^{*}$ 变化 $<2\%$ & 不敏感 \\
  上下文长度：$2{,}048\to131{,}072$ tokens（C7 可行取值） & 注意力/训练倍数 $0.07\to4.37$；$N^{*}$ 由 $6.089$B 降至 $2.705$B；损失上升 $6.0\%$ & 高度敏感，存在临界点 $L_{\text{crit}}=30{,}000$ \\
  算力预算 $C$：$10^{19}\to10^{24}$ FLOPs & $Q^{*}$ 由 $0.551$ 升至 $1.000$；$D^*/N^*$ 由 $30.9$ 升至 $76.1$；质量份额 $0\to0.407\to0.016$ & 高度敏感，发生三段式结构性转移 \\
  半合成数据口径：并入 B2 / B8 & B2 使 MAPE 升至 $33.7\%$；B8 使 MAPE 升至 $133.4\%$ 且预测低于不可约损失下限 & 不可并入；已划出可信度边界 \\
\end{longtable}

由表\ref{tab:sensitivity}可知，本文结论可按稳健程度分为三类：

\textbf{（1）高度稳健的结论：} 冲突率显著低于独立基准（任一阈值下比值均 $<0.93$）；岭正则系数不影响配比系数与拟合精度；质量基线口径不影响最优配置；中高预算下“质量拉满”的结论不依赖质量成本函数的具体形式。这些结论可直接用于后续决策。

\textbf{（2）存在临界点的结论：} 上下文长度跨越 $L_{\text{crit}}=30{,}000$ tokens 后注意力开销超过训练开销，最优规模被压缩 $2.25$ 倍；算力预算跨越 $1.78\times10^{19}$ 与 $7.08\times10^{19}$ 两个门槛时，最优质量分别由“基线不变”转为内点提升、再由内点解转为角点解。这些“临界/门槛”本身就是本文的核心发现，而非模型缺陷。

\textbf{（3）需要谨慎表述的结论：} 低预算档（$C=10^{19}$ FLOPs）的最优质量对质量成本函数形式高度敏感（$0.551\sim1.000$，即从“完全不投资质量”到“质量拉满”），因此该档的具体数值应作为区间而非点值使用；Loss—Benchmark 桥接的 $R^2$ 仅 $0.26$，其换算出的“损失改善等价于多少能力分”只具方向意义；领域质量排序对指标方向判定与赋权方案均为中等以上敏感，引用该排序时须同时引述其口径。

综上，除已明确标注的低预算档与桥接环节外，本文模型的资源配置与演进结论对参数扰动均保持稳定，具备可靠性与可解释性。
